% !TEX program = xelatex
\documentclass[11pt,a4paper]{article}
\usepackage{fontspec}
\setmainfont{Calibri}
\setmonofont[Scale=0.88]{Consolas}
\newfontfamily\symbolfont{Segoe UI Symbol}
\usepackage{polyglossia}\setdefaultlanguage{polish}
\usepackage[table]{xcolor}
\usepackage{booktabs,longtable,array,geometry,enumitem,graphicx,fancyhdr,caption}
\usepackage[most]{tcolorbox}
\PassOptionsToPackage{hyphens}{url}
\usepackage[hidelinks,unicode]{hyperref}
\emergencystretch=2.5em
\geometry{a4paper,margin=2.1cm,top=2.3cm,bottom=2.3cm}
\setlength{\parindent}{0pt}\setlength{\parskip}{5pt}
\definecolor{apteal}{HTML}{0E7C86}
\definecolor{apteallight}{HTML}{E6F4F5}
\definecolor{apgrey}{HTML}{F3F4F6}
\definecolor{apamber}{HTML}{B45309}
\definecolor{apamberlight}{HTML}{FEF3C7}
\definecolor{apred}{HTML}{B91C1C}
\definecolor{apredlight}{HTML}{FEE2E2}
\definecolor{apgreen}{HTML}{15803D}
\definecolor{apgreenlight}{HTML}{DCFCE7}
\renewcommand{\arraystretch}{1.25}
\captionsetup{font=small,labelfont=bf}
\pagestyle{fancy}\fancyhf{}
\fancyfoot[L]{\small AMPERE POINT — plan testów offline · Q11 z modułem Shelly · v1 · 2026-09-28}
\fancyfoot[R]{\small\thepage}
\renewcommand{\headrulewidth}{0pt}
\newtcolorbox{uwaga}[1][Uwaga]{enhanced,breakable,colback=apamberlight,colframe=apamber,title={#1},fonttitle=\bfseries,boxrule=0.6pt,arc=2pt,left=6pt,right=6pt,top=4pt,bottom=4pt}
\newtcolorbox{info}[1][Jak to działa]{enhanced,breakable,colback=apteallight,colframe=apteal,title={#1},fonttitle=\bfseries,boxrule=0.6pt,arc=2pt,left=6pt,right=6pt,top=4pt,bottom=4pt}
\newtcolorbox{wazne}[1][Ważne]{enhanced,breakable,colback=apredlight,colframe=apred,title={#1},fonttitle=\bfseries,boxrule=0.6pt,arc=2pt,left=6pt,right=6pt,top=4pt,bottom=4pt}
\newtcolorbox{przyklad}[1][Przykład liczbowy]{enhanced,breakable,colback=apgreenlight,colframe=apgreen,title={#1},fonttitle=\bfseries,boxrule=0.6pt,arc=2pt,left=6pt,right=6pt,top=4pt,bottom=4pt}
\newtcolorbox{kodbox}{enhanced,breakable,colback=apgrey,colframe=apgrey,boxrule=0pt,arc=2pt,left=5pt,right=5pt,top=3pt,bottom=3pt,fontupper=\ttfamily\footnotesize}
\setlist{nosep,leftmargin=*}
\setcounter{secnumdepth}{2}
\begin{document}

{\LARGE\bfseries Plan testów offline: Q11 z modułem Shelly\par}
{\large Co jeszcze trzeba sprawdzić, żeby potwierdzić, że bez internetu Q11 na Shelly umie wszystko to, co Q11 na Tuya, i że po zaniku zasilania pamięta to, co powinna · AMPERE POINT · 28 września 2026 · oprac. Dawid Cekała\par}
\vspace{4pt}\textcolor{apteal}{\rule{\textwidth}{1.2pt}}

\section*{W skrócie}

\begin{itemize}
\item \textbf{Co Tuya daje offline?} Prawie nic po stronie modułu: aplikacja Tuya bez chmury nie działa. Wszystko, co działa bez internetu, siedzi w sterowniku Q11: menu, przyciski, start po podłączeniu auta, limit prądu, okno godzin, limit energii, wyświetlacz. Moduł Tuya dokłada tylko ikonę sieci i godzinę.
\item \textbf{Co więc sprawdzamy?} Dwie rzeczy. Po pierwsze, że funkcje sterownika nadal działają, gdy w miejscu modułu Tuya siedzi Shelly, a internetu nie ma. Po drugie, że to, co moduł Shelly daje od siebie bez chmury, naprawdę działa: Bluetooth, harmonogramy w module, webhooki, pamięć ustawień.
\item \textbf{Trzecia sprawa: zanik zasilania.} Moduł jest zasilany z płytki, więc gaśnie razem z ładowarką. Po powrocie zasilania dwa układy startują naraz i każdy ma własną pamięć. Trzeba sprawdzić, kto komu narzuca wartości i czy nic nie ginie.
\item \textbf{Kolejność od bezpiecznych do ryzykownych.} Reset sieci z menu ładowarki zostawiamy na koniec, bo może skasować parowanie Wi-Fi modułu.
\end{itemize}

\section{Co ma Q11 na Tuya bez internetu, a co Q11 na Shelly}

\begin{longtable}{>{\raggedright\arraybackslash}p{4.56cm}>{\raggedright\arraybackslash}p{1.82cm}>{\raggedright\arraybackslash}p{2.74cm}>{\raggedright\arraybackslash}p{4.33cm}>{\raggedright\arraybackslash}p{0.95cm}}
\toprule
\rowcolor{apteallight}\textbf{Funkcja bez internetu} & \textbf{Gdzie siedzi} & \textbf{Q11 na Tuya} & \textbf{Q11 na Shelly} & \textbf{Test} \\
\midrule\endhead
\bottomrule\endfoot
Start ładowania po podłączeniu auta & sterownik & \leavevmode{\symbolfont\color{apgreen}\textbf{✔}} & \leavevmode{\symbolfont\color{apgreen}\textbf{✔}} sprawdzone 25.09 & T2 \\
Limit prądu z menu & sterownik & \leavevmode{\symbolfont\color{apgreen}\textbf{✔}} & ? nie sprawdzano z menu, tylko z modułu & T2 \\
Okno godzin (punkt 19) & sterownik & \leavevmode{\symbolfont\color{apgreen}\textbf{✔}} z menu; z aplikacji przez chmurę & ? & T3 \\
Limit energii sesji (punkt 17) & sterownik & \leavevmode{\symbolfont\color{apgreen}\textbf{✔}} & ? sterownik nigdy nie zgłosił punktu 17 & T3 \\
Tryb pracy: natychmiast, energia, harmonogram (punkt 14) & sterownik & \leavevmode{\symbolfont\color{apgreen}\textbf{✔}} & ? & T3 \\
Godzina dla sterownika & moduł & z chmury Tuya; bez internetu brak & \leavevmode{\symbolfont\color{apgreen}\textbf{✔}} z lokalnego serwera czasu, sprawdzone & T2 \\
Ikona sieci na wyświetlaczu & moduł & odbija stan chmury & \leavevmode{\symbolfont\color{apgreen}\textbf{✔}} moduł zawsze wysyła „4”; ikona świeci nawet bez internetu & T2 \\
Sterowanie z telefonu bez internetu & moduł & \leavevmode{\symbolfont\color{apred}\textbf{✘}} aplikacja Tuya tylko przez chmurę & ? Bluetooth w aplikacji Shelly; strona WWW modułu & T1 \\
Sterowanie z komputera w sieci lokalnej & moduł & \leavevmode{\symbolfont\color{apamber}\textbf{◐}} tylko nieoficjalnie, kluczem lokalnym & \leavevmode{\symbolfont\color{apgreen}\textbf{✔}} HTTP, WebSocket, MQTT, sprawdzone & — \\
Harmonogram w urządzeniu & — & \leavevmode{\symbolfont\color{apred}\textbf{✘}} tylko w chmurze & ? do 20 zadań w module & T4 \\
Powiadomienie o zdarzeniu do własnego systemu & — & \leavevmode{\symbolfont\color{apred}\textbf{✘}} & ? webhooki w module & T4 \\
Pamięć ustawień po zaniku zasilania & sterownik i moduł & sterownik pamięta swoje & ? dwa układy z pamięcią, do rozstrzygnięcia & T5 \\
Reset sieci z menu ładowarki & sterownik {\symbolfont→} moduł & \leavevmode{\symbolfont\color{apgreen}\textbf{✔}} moduł Tuya wchodzi w parowanie & ? co robi moduł Shelly na to polecenie & T6 \\
\end{longtable}

\section{Założenia i przygotowanie}

\begin{itemize}
\item Chmura Shelly wyłączona przez cały czas, jak dotąd.
\item W sieci działają kontenery: broker MQTT i lokalny serwer czasu (192.168.0.57). Serwer czasu musi działać \textbf{przed} każdym włączeniem ładowarki, bo moduł pyta o czas zaraz po starcie.
\item Monitor logu przez Wi-Fi uruchomiony przed każdym testem; do testów z zanikiem zasilania monitor sam się wznawia po powrocie modułu.
\item Symulator auta podłączony: stany „podłączone” (9 V) i „ładuje” (6 V). Bez obciążenia prądowego; prąd i moc faz zostają poza tym planem.
\item Laptop ma kartę Bluetooth (Intel); do testów z komputera biblioteka \texttt{bleak} w Pythonie. Telefon z aplikacją Shelly do testu z użytkownikiem.
\item Wi-Fi laptopa nie ruszamy. Zanik Wi-Fi symulujemy po stronie modułu, nie laptopa.
\item Po każdym teście zapis w tabeli wyników: \leavevmode{\symbolfont\color{apgreen}\textbf{✔}} działa, \leavevmode{\symbolfont\color{apamber}\textbf{◐}} działa z zastrzeżeniem, \leavevmode{\symbolfont\color{apred}\textbf{✘}} nie działa, z opisem tego, co zaobserwowano.
\end{itemize}

\section{Testy}

\subsection{T1. Bluetooth}

\textbf{Cel.} Potwierdzić, że ładowarką da się sterować bez Wi-Fi i bez internetu, z telefonu i z komputera.


\begin{longtable}{>{\raggedright\arraybackslash}p{3.00cm}>{\raggedright\arraybackslash}p{6.43cm}>{\raggedright\arraybackslash}p{5.85cm}}
\toprule
\rowcolor{apteallight}\textbf{Krok} & \textbf{Co robimy} & \textbf{Wynik oczekiwany} \\
\midrule\endhead
\bottomrule\endfoot
1a & Skan Bluetooth z laptopa & moduł widoczny pod swoją nazwą, z reklamą Shelly \\
1b & Polecenie przez Bluetooth z laptopa: odczyt stanu, potem zmiana limitu prądu na 10 A & sterownik potwierdza w logu w ciągu 1 s, tak jak przy Wi-Fi \\
1c & Telefon: Wi-Fi i dane komórkowe wyłączone, aplikacja Shelly przez Bluetooth: włącz/wyłącz ładowanie, zmiana limitu & to samo co 1b; sprawdzamy też, ile trwa połączenie \\
1d & Moduł bez Wi-Fi: przez Bluetooth wyłączamy stację Wi-Fi modułu (wcześniej upewniamy się, że własna sieć modułu, punkt dostępowy, jest włączona jako droga awaryjna); powtarzamy 1b i 1c & działa bez Wi-Fi; log niedostępny, więc dowodem jest wyświetlacz ładowarki i odpowiedź przez Bluetooth \\
1e & Przez Bluetooth włączamy stację Wi-Fi z powrotem & moduł wraca do sieci biura w ciągu minuty, log wraca \\
\end{longtable}

\textbf{Ryzyko.} Jeśli w 1d Bluetooth zawiedzie, moduł zostaje bez Wi-Fi. Droga powrotu: własna sieć modułu (punkt dostępowy) z telefonu. Dlatego 1d dopiero po udanych 1b i 1c.


\textbf{Narzędzie do napisania.} \texttt{narzedzia\textbackslash{}ble\_\allowbreak{}rpc.\allowbreak{}py}: klient poleceń Shelly przez Bluetooth (ta sama treść JSON co przez Wi-Fi, opakowana w ramkę Bluetooth Shelly).


\subsection{T2. Funkcje sterownika bez internetu}

\textbf{Cel.} Potwierdzić, że sterownik robi swoje, gdy moduł ma tylko sieć lokalną.


\begin{longtable}{>{\raggedright\arraybackslash}p{3.00cm}>{\raggedright\arraybackslash}p{4.66cm}>{\raggedright\arraybackslash}p{7.62cm}}
\toprule
\rowcolor{apteallight}\textbf{Krok} & \textbf{Co robimy} & \textbf{Wynik oczekiwany} \\
\midrule\endhead
\bottomrule\endfoot
2a & Symulator: wolne {\symbolfont→} podłączone {\symbolfont→} ładuje & sterownik przechodzi stany 0 {\symbolfont→} 1 {\symbolfont→} 4 (punkt 3), zamyka stycznik; moduł pokazuje to w polach \\
2b & Limit prądu z \textbf{menu ładowarki}: 8 A & sterownik zgłasza punkt 4 = 8 samoistnie; pole w module aktualizuje się bez naszego zapisu \\
2c & Limit prądu z modułu przy ładowaniu: 12 A & sterownik potwierdza w 1 s; wyświetlacz pokazuje 12 A \\
2d & Godzina po starcie: włączamy ładowarkę, patrzymy na pierwsze ramki czasu & do \textasciitilde{}20 s moduł odpowiada zerową datą, potem poprawną; sprawdzamy, co w tym czasie pokazuje wyświetlacz \\
2e & Ikona sieci: obserwacja & świeci przy samej sieci lokalnej; zapisujemy jako cechę do decyzji: czy „4” bez internetu to dobrze \\
\end{longtable}

\subsection{T3. Okno godzin, tryb pracy, limit energii}

\textbf{Cel.} Trzy funkcje, których Tuya obsługuje z aplikacji, a my na razie tylko z menu. Ustalić, czy sterownik je zgłasza i wykonuje, oraz co trzeba dodać w produkcie Shelly.


\begin{longtable}{>{\raggedright\arraybackslash}p{3.00cm}>{\raggedright\arraybackslash}p{6.14cm}>{\raggedright\arraybackslash}p{6.14cm}}
\toprule
\rowcolor{apteallight}\textbf{Krok} & \textbf{Co robimy} & \textbf{Wynik oczekiwany} \\
\midrule\endhead
\bottomrule\endfoot
3a & Z menu ładowarki: tryb „harmonogram”, okno np. od bieżącej godziny + 2 min do + 5 min; symulator w stanie „ładuje” & sterownik zgłasza punkt 14 = 2 i punkt 19 z godzinami; ładowanie rusza w oknie i staje po nim; poza oknem stan „czeka” \\
3b & To samo bez ważnego czasu: serwer czasu zatrzymany, ładowarka włączona na nowo & sterownik dostaje zerową datę; zapisujemy, co robi z oknem: nie startuje, startuje od razu, czy coś innego. Potem serwer czasu z powrotem \\
3c & Z menu: tryb „energia”, limit np. 1 kWh & sterownik zgłasza punkt 14 = 1 i, mamy nadzieję, punkt 17; bez obciążenia licznik nie ruszy, więc sprawdzamy tylko zgłaszanie \\
3d & Zapis punktów 14, 17 i 19 z modułu (surowym poleceniem do punktu danych, bez zmiany produktu) & sterownik przyjmuje i pokazuje w menu \\
\end{longtable}

\textbf{Wynik uboczny.} Lista pól do dodania w portalu: tryb pracy, okno godzin, limit energii, plus błędy, sygnał z auta i temperatura z poprzedniej listy.


\subsection{T4. Harmonogram, webhook i pamięć w module}

\textbf{Cel.} To, co moduł Shelly daje ponad Tuya, ma działać bez chmury i przetrwać restart.


\begin{longtable}{>{\raggedright\arraybackslash}p{3.00cm}>{\raggedright\arraybackslash}p{7.13cm}>{\raggedright\arraybackslash}p{5.15cm}}
\toprule
\rowcolor{apteallight}\textbf{Krok} & \textbf{Co robimy} & \textbf{Wynik oczekiwany} \\
\midrule\endhead
\bottomrule\endfoot
4a & Harmonogram w module: o +2 min ustaw limit 6 A, o +4 min 16 A & oba zadania odpalają o czasie, sterownik potwierdza \\
4b & Webhook: przy zmianie stanu ładowania moduł wywołuje adres na laptopie (mały nasłuch HTTP w Pythonie) & wywołanie przychodzi przy każdej zmianie, z treścią zdarzenia \\
4c & Pamięć klucz-wartość: zapis wartości testowej & odczyt daje to samo; sprawdzenie po T5, czy przetrwała zanik zasilania \\
4d & Harmonogram i webhook po restarcie modułu (samo polecenie restartu, bez zaniku zasilania) & nadal są i działają \\
\end{longtable}

\subsection{T5. Zanik zasilania}

\textbf{Cel.} Rozstrzygnąć, co pamięta sterownik, co moduł, i kto komu narzuca wartość po starcie.


\begin{longtable}{>{\raggedright\arraybackslash}p{3.00cm}>{\raggedright\arraybackslash}p{6.14cm}>{\raggedright\arraybackslash}p{6.14cm}}
\toprule
\rowcolor{apteallight}\textbf{Krok} & \textbf{Co robimy} & \textbf{Wynik oczekiwany / co zapisujemy} \\
\midrule\endhead
\bottomrule\endfoot
5a & Stan wyjściowy, zapisany z logu: limit 10 A ustawiony \textbf{z modułu}, ładowanie wyłączone, tryb „harmonogram” z oknem, limit energii 1 kWh, wpis w pamięci klucz-wartość, harmonogram z T4, symulator w stanie „podłączone” & pełny zrzut wszystkich punktów danych i pól modułu \\
5b & Wyłącznik ładowarki: wyłączyć na 30 s, włączyć & moduł i sterownik startują; monitor wraca \\
5c & Zrzut po starcie (odpowiedź sterownika na zapytanie o wszystkie punkty) porównany ze stanem 5a & które wartości sterownik zachował: oczekujemy limit prądu, tryb, okno, limit energii; zapisujemy każdą różnicę \\
5d & Pola modułu po starcie a wartości sterownika & pole „limit prądu” w module jest zapamiętywane w jego pamięci. Jeśli sterownik wstał z inną wartością niż moduł: kto wygrał? Czy moduł wysłał swoją wartość do sterownika (zapis w logu), czy przyjął jego? Zapis w logu rozstrzyga \\
5e & To samo dla włącznika ładowania: przed wyłączeniem „włączone”, symulator w stanie „ładuje” & czy po powrocie zasilania ładowanie rusza samo, czy czeka na polecenie; to jest zachowanie sterownika, ważne dla klienta \\
5f & Czas: ile sekund po starcie sterownik dostał poprawną godzinę & z logu; wcześniej zerowa data \\
5g & Moduł: konfiguracja Wi-Fi, MQTT, serwer czasu, harmonogram, webhook, wpis klucz-wartość & wszystko na miejscu \\
5h & Energia i czas sesji w polach modułu & oczekujemy zera: znany brak, do przeniesienia do pamięci klucz-wartość \\
5i & Powtórka 5a–5c z limitem ustawionym \textbf{z menu ładowarki} zamiast z modułu & sprawdzamy drugi kierunek: czy moduł po starcie nie nadpisze ustawienia z menu swoją zapamiętaną wartością \\
\end{longtable}

\textbf{Dlaczego to ważne.} Przykład: klient ustawia z menu 8 A, bo ma słabe zabezpieczenie. Moduł pamięta z aplikacji 16 A. Po zaniku zasilania moduł wstaje, wysyła „16 A” i ładowarka rusza z 16 A. Taki błąd byłby groźny, więc 5d i 5i to najważniejsze punkty planu.


\subsection{T6. Reset sieci z menu ładowarki}

\textbf{Cel.} W Q11 na Tuya pozycja menu „reset sieci” wprowadza moduł w parowanie. Sprawdzić, co polecenie sterownika robi z modułem Shelly.


\begin{longtable}{>{\raggedright\arraybackslash}p{3.00cm}>{\raggedright\arraybackslash}p{6.45cm}>{\raggedright\arraybackslash}p{5.83cm}}
\toprule
\rowcolor{apteallight}\textbf{Krok} & \textbf{Co robimy} & \textbf{Wynik oczekiwany} \\
\midrule\endhead
\bottomrule\endfoot
6a & Monitor na najwyższym poziomie logu; z menu ładowarki: reset sieci & w logu ramka od sterownika (kod resetu Wi-Fi) i reakcja modułu: ignoruje, kasuje Wi-Fi, wchodzi w parowanie \\
6b & Jeśli moduł skasował Wi-Fi: ponowne parowanie z telefonu przez Bluetooth, jak 24.09 & moduł wraca do sieci \\
\end{longtable}

\textbf{Ryzyko.} Utrata Wi-Fi modułu; dlatego na końcu. Co się stanie, zapisujemy jako fakt do dokumentacji dla fabryki (pytanie 8).


\section{Kolejność i czas}

\begin{longtable}{>{\raggedright\arraybackslash}p{3.00cm}>{\raggedright\arraybackslash}p{5.00cm}>{\raggedright\arraybackslash}p{1.34cm}>{\raggedright\arraybackslash}p{2.33cm}>{\raggedright\arraybackslash}p{2.74cm}}
\toprule
\rowcolor{apteallight}\textbf{Kolejność} & \textbf{Test} & \textbf{Ryzyko} & \textbf{Szacowany czas} & \textbf{Potrzebny użytkownik} \\
\midrule\endhead
\bottomrule\endfoot
1 & T1a–T1c Bluetooth przy działającym Wi-Fi & małe & 1 h z napisaniem narzędzia & tak, do 1c z telefonem \\
2 & T2 funkcje sterownika & małe & 30 min & tak, menu i symulator \\
3 & T5 zanik zasilania, pierwsza runda & małe & 45 min & tak, wyłącznik \\
4 & T3 okno, tryb, limit energii & małe & 45 min & tak, menu \\
5 & T4 harmonogram, webhook, pamięć & małe & 45 min & nie \\
6 & T5 druga runda, po T3 i T4 & małe & 30 min & tak \\
7 & T1d–T1e Bluetooth bez Wi-Fi & średnie & 30 min & tak, telefon jako droga awaryjna \\
8 & T6 reset sieci z menu & duże & 20 min plus ewentualne parowanie & tak \\
\end{longtable}

Razem około 5 godzin z przerwami. Po wszystkim: tabela wyników do analizy GAP (v4) i uzupełnienie dokumentacji dla fabryki o wyniki T5 i T6.


\section{Czego ten plan nie obejmuje}

\begin{itemize}
\item Prąd i moc faz pod obciążeniem, pełna sesja z energią: wymaga obciążenia, osobny test.
\item Aktualizacja sterownika przez moduł: wymaga pliku od fabryki.
\item Ekran ładowarki w aplikacji Shelly: sprawa portalu, nie offline.
\end{itemize}

\end{document}